% Options for packages loaded elsewhere
\PassOptionsToPackage{unicode}{hyperref}
\PassOptionsToPackage{hyphens}{url}
%
\documentclass[
  11pt,
  ignorenonframetext,
  aspectratio=169,
]{beamer}
\usepackage{pgfpages}
\setbeamertemplate{caption}[numbered]
\setbeamertemplate{caption label separator}{: }
\setbeamercolor{caption name}{fg=normal text.fg}
\beamertemplatenavigationsymbolsempty
% Prevent slide breaks in the middle of a paragraph
\widowpenalties 1 10000
\raggedbottom

\usepackage{amsmath,amssymb}
\usepackage{iftex}
\ifPDFTeX
  \usepackage[T1]{fontenc}
  \usepackage[utf8]{inputenc}
  \usepackage{textcomp} % provide euro and other symbols
\else % if luatex or xetex
  \usepackage{unicode-math}
  \defaultfontfeatures{Scale=MatchLowercase}
  \defaultfontfeatures[\rmfamily]{Ligatures=TeX,Scale=1}
\fi
\usepackage{lmodern}
\usefonttheme{serif}
\ifPDFTeX\else  
    % xetex/luatex font selection
\fi
% Use upquote if available, for straight quotes in verbatim environments
\IfFileExists{upquote.sty}{\usepackage{upquote}}{}
\IfFileExists{microtype.sty}{% use microtype if available
  \usepackage[]{microtype}
  \UseMicrotypeSet[protrusion]{basicmath} % disable protrusion for tt fonts
}{}
\makeatletter
\@ifundefined{KOMAClassName}{% if non-KOMA class
  \IfFileExists{parskip.sty}{%
    \usepackage{parskip}
  }{% else
    \setlength{\parindent}{0pt}
    \setlength{\parskip}{6pt plus 2pt minus 1pt}}
}{% if KOMA class
  \KOMAoptions{parskip=half}}
\makeatother
\usepackage{xcolor}
\newif\ifbibliography
\setlength{\emergencystretch}{3em} % prevent overfull lines
\setcounter{secnumdepth}{-\maxdimen} % remove section numbering


\providecommand{\tightlist}{%
  \setlength{\itemsep}{0pt}\setlength{\parskip}{0pt}}\usepackage{longtable,booktabs,array}
\usepackage{calc} % for calculating minipage widths
\usepackage{caption}
% Make caption package work with longtable
\makeatletter
\def\fnum@table{\tablename~\thetable}
\makeatother
\usepackage{graphicx}
\makeatletter
\def\maxwidth{\ifdim\Gin@nat@width>\linewidth\linewidth\else\Gin@nat@width\fi}
\def\maxheight{\ifdim\Gin@nat@height>\textheight\textheight\else\Gin@nat@height\fi}
\makeatother
% Scale images if necessary, so that they will not overflow the page
% margins by default, and it is still possible to overwrite the defaults
% using explicit options in \includegraphics[width, height, ...]{}
\setkeys{Gin}{width=\maxwidth,height=\maxheight,keepaspectratio}
% Set default figure placement to htbp
\makeatletter
\def\fps@figure{htbp}
\makeatother

\usepackage{tikz}
\usetikzlibrary{calc}

\definecolor{fondo}{HTML}{FAF6F1}
\definecolor{texto}{HTML}{2E2226}
\definecolor{vino}{HTML}{6F2C45}
\definecolor{champana}{HTML}{C4A27F}
\definecolor{beige}{HTML}{EFE4DA}
\definecolor{gris}{HTML}{7D7075}

\setbeamertemplate{navigation symbols}{}
\setbeamersize{text margin left=1.6cm,text margin right=1.2cm}

\setbeamercolor{normal text}{fg=texto,bg=fondo}
\setbeamercolor{structure}{fg=vino}
\setbeamercolor{frametitle}{fg=vino}
\setbeamercolor{footline}{fg=gris}
\setbeamercolor{block title}{fg=vino,bg=beige}
\setbeamercolor{block body}{fg=texto,bg=beige}
\setbeamertemplate{blocks}[default]
\setbeamertemplate{itemize item}{\color{champana}\rule[0.35ex]{0.9ex}{0.9ex}}
\setbeamertemplate{itemize subitem}{\color{champana}\textendash}

\setbeamertemplate{background canvas}{%
  \begin{tikzpicture}[remember picture,overlay]
    \fill[fondo] (current page.south west) rectangle (current page.north east);
    \fill[vino] (current page.south west) rectangle ($(current page.north west)+(0.55,0)$);
    \draw[champana,line width=0.6pt]
      ($(current page.south west)+(0.8,0)$) -- ($(current page.north west)+(0.8,0)$);
  \end{tikzpicture}}

\setbeamertemplate{frametitle}{%
  \nointerlineskip\vskip3.5ex
  \begin{beamercolorbox}[wd=\paperwidth,leftskip=1.6cm,rightskip=1.2cm]{frametitle}
    {\Large\bfseries\insertframetitle\par}
    \vskip1ex
    {\color{champana}\rule{1.2cm}{1.4pt}\color{beige}\rule{\dimexpr\paperwidth-4cm\relax}{0.5pt}}
  \end{beamercolorbox}}

\setbeamertemplate{footline}{%
  \ifnum\insertframenumber>1
    \leavevmode
    \hbox{%
      \begin{beamercolorbox}[wd=.7\paperwidth,ht=3ex,dp=1.5ex,leftskip=1.6cm]{footline}
      \end{beamercolorbox}%
      \begin{beamercolorbox}[wd=.3\paperwidth,ht=3ex,dp=1.5ex,rightskip=1.2cm]{footline}
        \hfill\scriptsize\insertframenumber
      \end{beamercolorbox}}%
    \vskip0.4cm
  \fi}

\setbeamertemplate{title page}{%
  \raggedright
  {\fontsize{24}{28}\selectfont\color{vino}\bfseries\inserttitle\par}
  \vskip1.5ex
  {\color{champana}\rule{3cm}{1.4pt}\par}
  \vskip1.5ex
  {\large\color{gris}\insertsubtitle\par}
  \vskip5ex
  {\color{texto}%
    Alexandra González Bermúdez\\
    Emily Estefanía Mora Contreras\\
    César Salazar\\
    Daniela Prado Vargas\par}}
\makeatletter
\@ifpackageloaded{caption}{}{\usepackage{caption}}
\AtBeginDocument{%
\ifdefined\contentsname
  \renewcommand*\contentsname{Tabla de contenidos}
\else
  \newcommand\contentsname{Tabla de contenidos}
\fi
\ifdefined\listfigurename
  \renewcommand*\listfigurename{Listado de Figuras}
\else
  \newcommand\listfigurename{Listado de Figuras}
\fi
\ifdefined\listtablename
  \renewcommand*\listtablename{Listado de Tablas}
\else
  \newcommand\listtablename{Listado de Tablas}
\fi
\ifdefined\figurename
  \renewcommand*\figurename{Figura}
\else
  \newcommand\figurename{Figura}
\fi
\ifdefined\tablename
  \renewcommand*\tablename{Tabla}
\else
  \newcommand\tablename{Tabla}
\fi
}
\@ifpackageloaded{float}{}{\usepackage{float}}
\floatstyle{ruled}
\@ifundefined{c@chapter}{\newfloat{codelisting}{h}{lop}}{\newfloat{codelisting}{h}{lop}[chapter]}
\floatname{codelisting}{Listado}
\newcommand*\listoflistings{\listof{codelisting}{Listado de Listados}}
\makeatother
\makeatletter
\makeatother
\makeatletter
\@ifpackageloaded{caption}{}{\usepackage{caption}}
\@ifpackageloaded{subcaption}{}{\usepackage{subcaption}}
\makeatother

\ifLuaTeX
\usepackage[bidi=basic]{babel}
\else
\usepackage[bidi=default]{babel}
\fi
\babelprovide[main,import]{spanish}
% get rid of language-specific shorthands (see #6817):
\let\LanguageShortHands\languageshorthands
\def\languageshorthands#1{}
\ifLuaTeX
  \usepackage{selnolig}  % disable illegal ligatures
\fi
\usepackage{bookmark}

\IfFileExists{xurl.sty}{\usepackage{xurl}}{} % add URL line breaks if available
\urlstyle{same} % disable monospaced font for URLs
\hypersetup{
  pdftitle={Aumento del beneficio por fallecimiento en un seguro colectivo},
  pdfauthor={Alexandra González Bermúdez; Emily Estefanía Mora Contreras; César Salazar Amaya; Daniela Prado Vargas},
  pdflang={es},
  hidelinks,
  pdfcreator={LaTeX via pandoc}}


\title{Aumento del beneficio por fallecimiento en un seguro colectivo}
\subtitle{NPV, efectivo y respaldo de las reservas}
\author{Alexandra González Bermúdez \and Emily Estefanía Mora
Contreras \and César Salazar Amaya \and Daniela Prado Vargas}
\date{}

\begin{document}
\frame{\titlepage}


\begin{frame}{Comparación del beneficio: 5 M frente a 5.5 M}
\phantomsection\label{comparaciuxf3n-del-beneficio-5-m-frente-a-5.5-m}
\footnotesize

\begin{longtable}[]{@{}
  >{\raggedright\arraybackslash}p{(\columnwidth - 6\tabcolsep) * \real{0.4915}}
  >{\raggedleft\arraybackslash}p{(\columnwidth - 6\tabcolsep) * \real{0.1695}}
  >{\raggedleft\arraybackslash}p{(\columnwidth - 6\tabcolsep) * \real{0.1695}}
  >{\raggedleft\arraybackslash}p{(\columnwidth - 6\tabcolsep) * \real{0.1695}}@{}}
\toprule\noalign{}
\begin{minipage}[b]{\linewidth}\raggedright
Indicador
\end{minipage} & \begin{minipage}[b]{\linewidth}\raggedleft
Beneficio 5.000.000
\end{minipage} & \begin{minipage}[b]{\linewidth}\raggedleft
Beneficio 5.500.000
\end{minipage} & \begin{minipage}[b]{\linewidth}\raggedleft
Diferencia
\end{minipage} \\
\midrule\noalign{}
\endhead
VP de beneficios (M) & 22 104,9280 & 24 315,4208 & 2 210,4928 \\
NPV (M) & -12 913,9692 & -15 124,4620 & -2 210,4928 \\
Margen & -0,5788 & -0,6779 & -0,0991 \\
Saldo mínimo antes de aportes (M) & -1 591,9792 & -2 099,0550 &
-507,0758 \\
Total de aportes (M) & 3 051,3687 & 5 610,5287 & 2 559,1600 \\
Respaldo año 1 (M) & -2 024,8493 & -2 510,1508 & -485,3015 \\
\bottomrule\noalign{}
\end{longtable}

\normalsize

M = millones de unidades monetarias. El NPV baja exactamente lo que sube
el VP de beneficios.
\end{frame}

\begin{frame}{¿Es defendible con la prima de 50 000?}
\phantomsection\label{es-defendible-con-la-prima-de-50-000}
\footnotesize

\begin{longtable}[]{@{}
  >{\raggedright\arraybackslash}p{(\columnwidth - 4\tabcolsep) * \real{0.6562}}
  >{\raggedleft\arraybackslash}p{(\columnwidth - 4\tabcolsep) * \real{0.1719}}
  >{\raggedleft\arraybackslash}p{(\columnwidth - 4\tabcolsep) * \real{0.1719}}@{}}
\toprule\noalign{}
\begin{minipage}[b]{\linewidth}\raggedright
Indicador
\end{minipage} & \begin{minipage}[b]{\linewidth}\raggedleft
Base
\end{minipage} & \begin{minipage}[b]{\linewidth}\raggedleft
Mayor
\end{minipage} \\
\midrule\noalign{}
\endhead
Económico: NPV (M) & -12 913,9692 & -15 124,4620 \\
Efectivo: primer aporte, año 3 (M) & 57,9189 & 1 607,1994 \\
Efectivo: total de aportes (M) & 3 051,3687 & 5 610,5287 \\
Respaldo: recursos / reserva al cierre del año 1 & 49,34 \% & 37,20
\% \\
Costo esperado de fallecimientos por contrato, año 1 & 46 661,3602 & 51
327,4962 \\
\bottomrule\noalign{}
\end{longtable}

\normalsize

\begin{itemize}
\tightlist
\item
  La prima neta de gasto y comisión (37 500 por contrato) no cubre el
  costo esperado ni en el caso base; con el beneficio mayor, el costo
  supera incluso la prima completa de 50 000.
\item
  Las tres medidas (NPV, aportes y respaldo) son distintas y las tres
  empeoran.
\end{itemize}
\end{frame}

\begin{frame}{Decisión: posponer el aumento}
\phantomsection\label{decisiuxf3n-posponer-el-aumento}
\small

\begin{itemize}
\tightlist
\item
  \textbf{Decisión.} Posponer el aumento mientras la prima siga en 50
  000.
\item
  \textbf{Condición.} Reconsiderarlo cuando el esquema con beneficio de
  5 500 000 tenga un NPV no negativo. El aumento resta \(2\,210{,}4928\)
  M, así que el caso base necesitaría un NPV de al menos esa cifra; hoy
  es de \(-12\,913{,}9692\) M.
\item
  \textbf{Indicador de seguimiento.} Respaldo del año 1:
  \(-2\,024{,}8493\) M en el base y \(-2\,510{,}1508\) M con el aumento.
\item
  \textbf{Operación verificable.} Año 1, beneficio mayor: recursos antes
  del aporte menos reserva
  \(\approx 1\,487{,}0119 - 3\,997{,}1626 = -2\,510{,}1508\) M.
\item
  \textbf{Limitación.} Los supuestos y las reservas son los del
  ejercicio: es una comparación, no una medida de suficiencia actuarial
  ni de solvencia regulatoria.
\end{itemize}
\end{frame}




\end{document}
